
\documentclass{article}
\usepackage{colortbl}
\usepackage{makecell}
\usepackage{multirow}
\usepackage{supertabular}

\begin{document}

\newcounter{utterance}

\centering \large Interaction Transcript for game `cladder', experiment `full\_v1.5\_default', episode 1876 with llama{-}3.1{-}8b{-}instruct.
\vspace{24pt}

{ \footnotesize  \setcounter{utterance}{1}
\setlength{\tabcolsep}{0pt}
\begin{supertabular}{c@{$\;$}|p{.15\linewidth}@{}p{.15\linewidth}p{.15\linewidth}p{.15\linewidth}p{.15\linewidth}p{.15\linewidth}}
    \# & \multicolumn{2}{c}{Player} && \multicolumn{2}{c}{Game Master} \\
    \hline

    \theutterance \stepcounter{utterance}  
    & & & \multicolumn{4}{p{0.6\linewidth}}{
        \cellcolor[rgb]{0.9,0.9,0.9}{
            \makecell[{{p{\linewidth}}}]{
                \texttt{\tiny{[P1$\langle$GM]}}
                \texttt{You are an expert in causal inference. The following question is not a typical commonsense query, but rather a meticulously designed question created by a professor specializing in causal inference, intended to assess the students' mastery of the course content.} \\
\texttt{Imagine a self{-}contained, hypothetical world with only the following conditions, and without any unmentioned factors or causal relationships: Medication has a direct effect on blood pressure and heart condition. Blood pressure has a direct effect on heart condition. For people not taking any medication, the probability of healthy heart is 64\%. For people taking medication, the probability of healthy heart is 21\%. For people taking medication, would it be more likely to see healthy heart if the person had taken no medication?} \\
\texttt{Start your answer with 'Yes' or 'No', followed by additional reasoning or evidence to support your explanation.} \\
            }
        }
    }
     \\ \\

    \theutterance \stepcounter{utterance}  
    & \multicolumn{4}{p{0.6\linewidth}}{
        \cellcolor[rgb]{0.561,0.737,0.561}{
            \makecell[{{p{\linewidth}}}]{
                \texttt{\tiny{[P1$\rangle$GM]}}
                \texttt{No.} \\
\\ 
\texttt{To answer this question, we need to consider the causal relationships and the conditional probabilities given. We are given that:} \\
\\ 
\texttt{1. Medication has a direct effect on blood pressure and heart condition.} \\
\texttt{2. Blood pressure has a direct effect on heart condition.} \\
\texttt{3. For people not taking any medication, the probability of healthy heart is 64\%.} \\
\texttt{4. For people taking medication, the probability of healthy heart is 21\%.} \\
\\ 
\texttt{We want to know whether it is more likely to see a healthy heart if the person had taken no medication, given that they are taking medication. This is a counterfactual question, and we can use the potential outcomes framework to answer it.} \\
\\ 
\texttt{Let's denote the potential outcomes as follows:} \\
\\ 
\texttt{{-} H: Healthy heart} \\
\texttt{{-} UH: Unhealthy heart} \\
\texttt{{-} M: Taking medication} \\
\texttt{{-} NM: Not taking medication} \\
\\ 
\texttt{We are given the following probabilities:} \\
\\ 
\texttt{P(H|NM) = 0.64 (probability of healthy heart given no medication)} \\
\texttt{P(H|M) = 0.21 (probability of healthy heart given medication)} \\
\\ 
\texttt{We want to know P(H|NM,M), which is the probability of healthy heart given that the person had taken no medication and is now taking medication.} \\
\\ 
\texttt{Using the law of total probability, we can write:} \\
\\ 
\texttt{P(H|NM,M) = P(H|NM) \textbackslash{}* P(NM|M) + P(H|M) \textbackslash{}* P(M|M)} \\
\\ 
\texttt{Since the person is taking medication, P(NM|M) = 0. However, we can rewrite the equation as:} \\
\\ 
\texttt{P(H|NM,M) = P(H|M) \textbackslash{}* P(M|M) / P(M|M)} \\
\\ 
\texttt{Since P(M|M) = 1 (the person is taking medication), we have:} \\
\\ 
\texttt{P(H|NM,M) = P(H|M) = 0.21} \\
\\ 
\texttt{This means that, given that the person is taking medication, it is not more likely to see a healthy heart if the person had taken no medication. The probability of a healthy heart is the same as the probability of a healthy heart given that the person is taking medication, which is 21\%.} \\
            }
        }
    }
    & & \\ \\

    \theutterance \stepcounter{utterance}  
    & & & \multicolumn{2}{p{0.3\linewidth}}{
        \cellcolor[rgb]{0.9,0.9,0.9}{
            \makecell[{{p{\linewidth}}}]{
                \texttt{\tiny{[GM$|$GM]}}
                \texttt{No} \\
            }
        }
    }
    & & \\ \\

    \theutterance \stepcounter{utterance}  
    & & & \multicolumn{2}{p{0.3\linewidth}}{
        \cellcolor[rgb]{0.9,0.9,0.9}{
            \makecell[{{p{\linewidth}}}]{
                \texttt{\tiny{[GM$|$GM]}}
                \texttt{yes} \\
            }
        }
    }
    & & \\ \\

    \theutterance \stepcounter{utterance}  
    & & & \multicolumn{2}{p{0.3\linewidth}}{
        \cellcolor[rgb]{0.9,0.9,0.9}{
            \makecell[{{p{\linewidth}}}]{
                \texttt{\tiny{[GM$|$GM]}}
                \texttt{game\_result = LOSE} \\
            }
        }
    }
    & & \\ \\

\end{supertabular}
}

\end{document}
